%% LyX 2.5.1 created this file.  For more info, see https://www.lyx.org/.
%% Do not edit unless you really know what you are doing.
\documentclass[english,hebrew]{article}
\usepackage[HE8,T1]{fontenc}
\usepackage[utf8]{inputenc}
\usepackage{babel}
\usepackage{amssymb}
\usepackage[unicode=true]
 {hyperref}
\begin{document}
\title{מבוא לטופולוגיה קבוצתית: טופולוגיות תת-המרחב וטופולוגיית המכפלה}
\maketitle
\begin{description}
\item [{קטגוריות:}] טופולוגיה
\item [{תגים:}] מרחבים טופולוגיים, מרחבי מכפלה
\item [{מזהה:}] \L{topological\_spaces\_subspace\_and\_product\_topology}
\end{description}
בפוסטים האחרונים התחלנו לדבר על מרחבים טופולוגיים בדרך המקובלת במתמטיקה:
נתנו הגדרה )\textquotedblleft קבוצות פתוחות\textquotedblright (, ראינו
כמה דוגמאות )\textquotedblleft מרחבים מטריים\textquotedblright , \textquotedblleft טופולוגיית
הסדר\textquotedblright ( ודיברנו על פונקציות שמשמרות מבנה )\textquotedblleft הומיאומורפיזמים\textquotedblright (.
מה עכשיו, מה עכשיו, מה עכשיו? אה, כמובן, איך שכחנו את החלק הכי כיפי
- איך בונים מרחבים חדשים מתוך מרחבים קיימים.

למי שמכירים תורת הקבוצות, מה שנעשה פה לא הולך להיות מפתיע כי נשתמש
בפעולות סטנדרטיות על קבוצות: ניקח \textbf{תת-קבוצות} וניקח \textbf{מכפלות}
של קבוצות )במקור תכננתי גם לקחת \textbf{מנות} אבל הפוסט נהיה ארוך
מדי גם ככה וזה נושא שאפשר לדבר עליו בנפרד(. מה שהופך את זה למעניין
הוא שעכשיו אנחנו צריכים למצוא דרך להגדיר \textbf{טופולוגיה} על הקבוצות
שנקבל - אם כי הדרך הזו תהיה טבעית למדי )עבור מכפלות נראה שיש שתי דרכים
ודווקא זו שנראית פחות טבעית היא השימושית והנפוצה(. ומן הסתם, כדאי
שיהיו לנו דוגמאות מול העיניים.

בואו נתחיל.

\section{טופולוגיית תת-המרחב}

נתחיל עם הדבר המתבקש: יש לנו מרחב טופולוגי \L{$\left(X,\mathcal{T}_{X}\right)$}
ויש לנו תת-קבוצה \L{$Y\subseteq X$}. איך נהפוך את \L{$Y$} למרחב
טופולוגי בפני עצמו, כלומר איך נגדיר \L{$\mathcal{T}_{Y}$} בצורה \textquotedblleft טבעית\textquotedblright{}
שתשמור על הדמיון ל-\L{$X$}? ובכן, זה קל:
\begin{itemize}
\item \L{$\mathcal{T}_{Y}=\left\{ U\cap Y\ |\ U\in\mathcal{T}_{X}\right\} $}
\end{itemize}
לא קשה להראות שאם יש לנו בסיס \L{$\mathcal{B}$} של \L{$X$}, אז \L{$\left\{ B\cap Y\ |\ B\in\mathcal{B}\right\} $}
הוא בסיס של \L{$Y$} עם טופולוגיית תת-המרחב, כלומר קל לנו גם לדעת
איך בסיס ל-\L{$Y$} נראה אם אנחנו מכירים בסיס ל-\L{$X$}.

הנה דוגמא פשוטה שכבר דיברנו עליה כשראינו את טופולוגיית הסדר. ניקח
את \L{$X=\mathbb{R}$} עם הטופולוגיה הרגילה, כלומר זו שהקטעים הפתוחים
הם בסיס עבורה, וניקח את תת-המרחב \L{$Y=\mathbb{R}^{\ge0}$}, כלומר
כל המספרים הממשיים האי-שליליים. אז עכשיו למשל הקבוצה הפתוחה \L{$\left(-1,1\right)$},
כשחותכים אותה עם \L{$Y$}, נותנת את \L{$[0,-1)$}, כלומר את הקטע החצי
פתוח והחצי \textbf{סגור}. זה תואם את מה שראינו עם טופולוגיית הסדר,
שבה כשהיה איבר מינימלי היינו צריכים לקחת בתור קבוצות פתוחות גם את
הקטעים הסגורים שמתחילים בו. באופן כללי, הבסיס לטופולוגייה הזו יכלול
את כל הקטעים מהצורה \L{$[0,a)$} עבור \L{$a>0$} ואת כל הקטעים מהצורה
\L{$\left(a,b\right)$} עבור \L{$0<a<b$} - בדיוק כמו טופולוגיית הסדר.

יש כאן משהו טיפה טריקי - \L{$[0,-1)$} היא קבוצה פתוחה ב-\L{$Y$}
אבל היא \textbf{לא} קבוצה פתוחה ב-\L{$X$} )כי ב-\L{$X$} אין ל-{\beginL 0\endL}
סביבה שכולה מוכלת ב-\L{$X$}(. אז צריך להיות זהירים בסיטואציות עם
תתי-מרחבים וכשאומרים שקבוצה היא \textquotedblleft פתוחה\textquotedblright{}
צריך להסביר איפה. מה שכן, קל לנחש איזה תנאי צריך לדרוש על \L{$Y$}
כדי להבטיח שכל קבוצה פתוחה ב-\L{$Y$} תהיה פתוחה גם ב-\L{$X$} - אם
\L{$Y$} עצמו פתוחה ב-\L{$X$} אז מן הסתם גם החיתוך \L{$U\cap Y$}
יהיה פתוח ב-\L{$X$} לכל \L{$U$} שפתוחה ב-\L{$X$} )כי חיתוך שתי
קבוצות פתוחות הוא קבוצה פתוחה(.

עד כאן הכל פשוט וברור, נכון? אבל זו טופולוגיה. בטופולוגיה הכל משוגע.
הדוגמא שנתתי למעלה היא \textbf{פשוטה מדי}. היא נותנת את הרושם שאם
יש לי קבוצה סדורה \L{$X$} עם טופולוגיית הסדר עליה, ואני מסתכל על
תת-קבוצה \L{$Y$}, אז אני אקבל את אותה טופולוגיה עליה בשתי דרכים שונות:
\begin{itemize}
\item אני אגדיר על \L{$Y$} את טופולוגיית תת-המרחב בהתאם לטופולוגיית הסדר
שיש על \L{$X$}.
\item אני אגדיר על \L{$Y$} את טופולוגיית הסדר בהתאם ליחס הסדר שיש על \L{$X$}.
\end{itemize}
בגישה הראשונה מתחילים מיחס סדר על \L{$X$}; הוא מייצר טופולוגיה על
\L{$X$}; ואז \L{$X$} \textquotedblleft מוריש\textquotedblright{}
את \textbf{הטופולוגיה} שלו ל-\L{$Y$}. בגישה השניה גם מתחילים מיחס
סדר על \L{$X$}, ואז \L{$X$} \textquotedblleft מוריש\textquotedblright{}
את \textbf{יחס הסדר} שלו ל-\L{$Y$}; ואז יחס הסדר הנורש מייצר טופולוגיה
על \L{$Y$}.

באופן כללי שתי השיטות הללו יכולות \textbf{לתת דברים שונים}.

בואו נראה דוגמא. לא צריך ללכת רחוק: אפשר להתחיל עם \L{$X=\mathbb{R}$}
אבל בתור \L{$Y$} להצטמצם קצת יותר מאשר סתם לחתוך את \L{$X$} באמצע.
ספציפית, בואו ניקח גם \textbf{סינגלטון}. נגדיר

\L{$Y=[0,1)\cup\left\{ 2\right\} $}

עכשיו, בטופולוגיית תת-המרחב, מכיוון ש-\L{$\left(1,3\right)\cap Y=\left\{ 2\right\} $}
ו-\L{$\left(1,3\right)$} פתוחה ב-\L{$X$}, נקבל ש-\L{$\left\{ 2\right\} $}
פתוחה ב-\L{$Y$}. מצד שני, \L{$Y$} יורש את יחס הסדר \textquotedblleft הרגיל\textquotedblright{}
של \L{$X$}. על פי טופולוגיית הסדר, הקבוצות הפתוחות של \L{$Y$} הן
כל הקטעים מהצורה \L{$\left(a,b\right)$} כך ש-\L{$a<b\in Y$} ובנוסף
לכך, מכיוון ש-\L{$2$} הוא האיבר הגדול ביותר, כל הקטעים מהצורה \L{$(a,2]$}
עבור \L{$a\in Y$}. למשל, \L{$(\frac{1}{2},2]=\left(\frac{1}{2},1\right)\cup\left\{ 2\right\} $}
הולך להיות קבוצה פתוחה, אבל \L{$\left\{ 2\right\} $} עצמו? לא. אין
שום דרך לקבל אותו בתור \L{$(a,2]$}. מה שחמוד כאן במיוחד הוא שאם תת-המרחב
שלנו היה \L{$Y=\left[0,1\right]\cup\left\{ 2\right\} $}, אז דווקא
היה אפשר לקחת את \L{$(1,2]=\left\{ 2\right\} $} בתור איבר בסיס; אבל
מכיוון ש-\L{$Y=[0,1)\cup\left\{ 2\right\} $}, אז \L{$1$} הוא בכלל
לא חלק מהמרחב ואי אפשר לקחת קטע שמתחיל בו.

אוקיי, זה מוזר. האם אפשר למצוא קריטריון שיבטיח לנו שהמוזר הזה לא מתרחש?
כן, זה לא קשה במיוחד. בדוגמא הנגדית שלנו הבעיה הייתה הניתוק הזה בין
{\beginL 2\endL} ובין שאר המרחב. כדי למנוע מזה להתרחש, אפשר לדרוש
שהמרחב יהיה \textbf{קבוצה קמורה} - הרעיון בקבוצה כזו באופן כללי הוא
שעבור כל שתי נקודות, הקו שמחבר אותן עובר כולו בתוך הקבוצה. בהקשר שלנו,
של קבוצות סדורות לינארית, תת-הקבוצה \L{$Y\subseteq X$} היא \textbf{קמורה}
ב-\L{$X$} אם לכל \L{$a,b\in Y$} מתקיים \L{$\left(a,b\right)\subseteq Y$},
כאשר \L{$\left(a,b\right)=\left\{ x\in X\ |\ a<x<y\right\} $} )כלומר
- אנחנו מסתכלים על כל הנקודות של המרחב הגדול \L{$X$}, לא רק על אלו
ששייכות ל-\L{$Y$}(.

בואו נוכיח שאם \L{$Y$} קמורה ב-\L{$X$}, אז שתי הטופולוגיות שאפשר
לקבל בדרכים שתיארתי קודם הן זהות. איך מראים שטופולוגיות הן זהות? אפשר
לקחת קבוצה פתוחה כללית בטופולוגיה אחת ולהראות שהיא פתוחה בטופולוגיה
השניה, אבל זה יכול להיות קשה כי קבוצה פתוחה כללית זה משהו מסובך. בשביל
זה יש לנו בסיסים ותת-בסיסים, כדי לפשט הוכחות כאלו. אם יש לנו שתי טופולוגיות
\L{$\mathcal{T}_{1},\mathcal{T}_{2}$} ויש לנו איבר \L{$U\in\mathcal{T}_{1}$}
ששייך לתת-בסיס לטופולוגיה \L{$\mathcal{T}_{1}$} ואנחנו מראים ש-\L{$U\in\mathcal{T}_{2}$},
סיימנו; כי כל איבר ב-\L{$\mathcal{T}_{1}$} אפשר לכתוב בתור חיתוכים
סופיים ואיחודים של איבר תת-בסיס של \L{$\mathcal{T}_{1}$}, ואם כל
האיברים הללו הם פתוחים ב-\L{$\mathcal{T}_{2}$} גם כל חיתוך סופי ואיחוד
שלהם יהיה כזה.

עבור טופולוגיית הסדר, תת-הבסיס בנוי על \textbf{קרניים}, קבוצות מהצורה
\L{$\left(a,\infty\right)$} ו-\L{$\left(-\infty,b\right)$}. מספיק
להראות ש-\L{$\left(a,\infty\right)$} פתוחה גם בטופולוגיית תת-המרחב,
כי ההוכחה עבור \L{$\left(-\infty,b\right)$} תהיה סימטרית. 

עכשיו, בואו נסתכל על זה בזהירות. אני לקחתי את הקרן \L{$\left(a,\infty\right)$},
אבל \textbf{באיזה מרחב}? אני כרגע במרחב \L{$Y$} שהוא תת-קבוצה של
\L{$X$} עם יחס הסדר שירשנו מ-\L{$X$}. כלומר הקרן היא

\L{$\left(a,\infty\right)=\left\{ y\in Y\ |\ a<y<\infty\right\} $}

עכשיו, כל איבר של \L{$Y$} הוא גם איבר של \L{$X$} אז אפשר לכתוב

\L{$\left(a,\infty\right)=\left\{ x\in X\ |\ a<x<\infty\right\} \cap Y$}

אבל הקבוצה \L{$\left\{ x\in X\ |\ a<x<\infty\right\} $} היא קרן במרחב
\L{$X$}, ולכן אם לוקחים את טופולוגיית הסדר על \L{$X$}, הקבוצה הזו
תהיה פתוחה ב-\L{$X$}, ולכן החיתוך שלה עם \L{$Y$} יהיה פתוח ב-\L{$Y$}
\textbf{בטופולוגיית תת-המרחב}, שזה מה שרצינו להראות.

עכשיו נלך על הכיוון השני - ניקח את \L{$Y$} עם טופולוגיית תת המרחב.
האם יש לנו תת-בסיס עבור הטופולוגיה הזו? בערך: אנחנו יודעים שכל קבוצה
פתוחה ב-\L{$Y$} מתקבלת מחיתוך של \L{$Y$} עם קבוצה פתוחה ב-\L{$X$},
ועל \L{$X$} יש לנו תת-בסיס כי יש עליו את טופולוגיית הסדר. לכן, אם
נחתוך את האיברים של כל תת-הבסיס הזה עם \L{$Y$} נקבל תת-בסיס של \L{$Y$}
)כי כל קבוצה פתוחה ב-\L{$Y$} מתקבלת מכך שעושים איחודים וחיתוכים סופיים
של אברי תת הבסיס של \L{$X$} \textbf{ובסוף} חותכים עם \L{$Y$}; אפשר
גם לחתוך מראש עם \L{$Y$} את כל הקבוצות המעורבות וזה יעבוד באותה מידה(.
תת-הבסיס של \L{$X$} כולל קרניים, למשל מהצורה \L{$\left(a,\infty\right)$},
ולכן תת-הבסיס של \L{$Y$} יכיל קבוצות מהצורה \L{$\left(a,\infty\right)\cap Y$}.
מה אנחנו יודעים על הקבוצות הללו?

ראשית, ייתכן כמובן ש-\L{$\left(a,\infty\right)\cap Y=\emptyset$},
וזה בסדר כי זו קבוצה פתוחה. ייתכן גם ש-\L{$\left(a,\infty\right)\cap Y=\left(a,\infty\right)$},
כלומר לא איבדנו איברים בחיתוך עם \L{$Y$}, וגם זה בסדר כי קרניים ב-\L{$Y$}
הן פתוחות בטופולוגיית הסדר. הסיטואציה הבעייתית היא אם אחרי החיתוך
אנחנו מקבלים משהו שהוא לא קרן ב-\L{$Y$}. זו הדוגמא שראינו קודם: היה
לנו את \L{$Y=[0,1)\cup\left\{ 2\right\} $} ואם למשל נחתוך אותו עם
\L{$\left(1,\infty\right)$} נקבל את \L{$\left\{ 2\right\} $} שהיא
לא קרן ב-\L{$Y$}. זה משהו שבלבל גם אותי - מה זאת אומרת זו לא קרן
ב-\L{$Y$}? למה אי אפשר להגיד שזו הקרן \L{$\left(1,\infty\right)$}
ב-\L{$Y$}? אה! כי \L{$1\notin Y$} אז אי אפשר להשתמש בו בתור נקודת
קצה. תזכרו, קרן ב-\L{$Y$} היא קבוצה מהצורה \L{$\left\{ y\in Y\ |\ y>a\right\} $}
עבור \L{$a\in Y$}, וזה בדיוק מה שלא קורה כאן. אני יכול, כמובן, להסתכל
על קרן כמו \L{$\left(\frac{1}{2},\infty\right)$}, אבל הבעיה היא שהקרן
הזו תכיל \textbf{עוד} נקודות חוץ מאשר את {\beginL 2\endL} עצמה ולכן
גם פה לא אקבל את \L{$\left\{ 2\right\} $}. אז זו הבעיה - החיתוך \L{$\left(a,\infty\right)\cap Y$}
עלול במקרה הכללי לא להיות קרן, וצריך להשתמש בקמירות של \L{$Y$} כדי
להראות שהוא כן קרן.

אז מהי \L{$\left(a,\infty\right)\cap Y$} אם \L{$Y$} קמורה? אם \L{$a\in Y$}
אין שום בעיה, עדיין נקבל קרן: \L{$\left(a,\infty\right)\cap Y=\left\{ y\in Y\ |\ y>a\right\} $}.
זו קרן גם ב-\L{$Y$} בזכות זה ש-\L{$a\in Y$}. הבעיה היא אם \L{$a\notin Y$},
אבל במקרה הזה \L{$a$} לא יכול להיות ב\textquotedblright חור באמצע\textquotedblright{}
של \L{$Y$} בגלל הקמירות של \L{$Y$}. כלומר, לא ייתכן שיש \L{$b,c\in Y$}
כך ש-\L{$b<a<c$} כי אז היינו מקבלים ש-\L{$a\in Y$} מהקמירות של \L{$Y$}.
אז יש בדיוק שתי אפשרויות: או שלכל \L{$b\in Y$} מתקיים \L{$a<b$},
או שלכל \L{$b\in Y$} מתקיים \L{$b<a$}. במקרה השני, \L{$\left(a,\infty\right)$}
לא כוללת בכלל איברים של \L{$Y$} )כי כל מה שיש שם גדול מ-\L{$a$}
שגדול מכל אברי \L{$Y$}( ולכן \L{$\left(a,\infty\right)\cap Y=\emptyset$}
וזה כאמור בסדר. במקרה הראשון, \L{$a$} הוא חסם מלמטה של \textbf{כל}
\L{$Y$} ולכן \L{$\left(a,\infty\right)$} מכיל את \textbf{כל} \L{$Y$}
כך ש-\L{$\left(a,\infty\right)\cap Y=Y$} וגם זה בסדר כי גם כל המרחב
הוא קבוצה פתוחה. זה מסיים את ההוכחה הזו.

\section{טופולוגיית המכפלה - החלק הקל}

אחת הבניות הבסיסיות של קבוצה חדשה מקבוצות קיימות היא \textbf{מכפלה
קרטזית}. אם \L{$A,B$} הן קבוצות כלשהן אז המכפלה הקרטזית שלהן היא
\L{$A\times B=\left\{ \left(a,b\right)\ |\ a\in A,b\in B\right\} $}
כאשר \L{$\left(a,b\right)$} הוא \textbf{זוג סדור} של איברים. על ידי
חזרה כמה פעמים על אותו טריק, אפשר לבנות מכפלות קרטזיות של מספר סופי
כלשהו של קבוצות, \L{$A_{1}\times A_{2}\times\ldots\times A_{n}$}
הוא אוסף כל ה-\L{$n$}-יות הסדורות \L{$\left(a_{1},\ldots,a_{n}\right)$}
כך ש-\L{$a_{i}\in A_{i}$}. הבניה הזו היא מאוד שימושית - באלגברה לינארית
רואים בערך מהרגע הראשון איך בונים את \L{$\mathbb{R}^{n}$}, המרחב
הממשי ה-\L{$n$}-ממדי שמכליל את הישר הממשי והמישור הממשי. בלינארית
מגדירים \textbf{מבנה אלגברי} על הקבוצה הזו, כלומר פעולות של חיבור
וכפל בסקלר; אנחנו רוצים להגדיר \textbf{טופולוגיה} על הקבוצות הללו,
בהנחה שכבר יש טופולוגיה על כל קבוצה בנפרד. איך נעשה את זה? ובכן, זה
אחד מהמקרים שבהם ההגדרה ה\textquotedblright טבעית\textquotedblright{}
היא גם זו המתבקשת - אבל רק כל עוד אנחנו \textbf{במקרה הסופי}. בקרוב
דברים יתחילו להסתבך.

אם כן, בואו נתחיל ממה שקל. נניח שיש לנו שני מרחבים טופולוגיים \L{$\left(X,\mathcal{T}_{X}\right),\left(Y,\mathcal{T}_{Y}\right)$}.
אנחנו לוקחים את המכפלה \L{$X\times Y$} ומגדירים עליה טופולוגיה, \L{$\mathcal{T}$}
על ידי זה שניקח בתור בסיס \L{$\mathcal{B}$} לטופולוגיה את
\begin{itemize}
\item \L{$\mathcal{B}=\left\{ U\times V\ |\ U\in\mathcal{T}_{X},V\in\mathcal{T}_{Y}\right\} $}
\end{itemize}
כלומר, הקבוצות הפתוחות הבסיסיות שלנו יהיו מכפלות של קבוצות פתוחות
משני המרחבים. שימו לב שזה רק \textbf{בסיס}, זה לחלוטין לא תיאור של
כל הקבוצות בטופולגיה; בכל מכפלה קרטזית יש הרבה מאוד תת-קבוצות שאי
אפשר לתאר בתור מכפלה של קבוצה מפה וקבוצה משם )למשל, במכפלה \L{$A^{2}$}
אין דרך לתאר את הקבוצה \L{$\left\{ \left(a,a\right)\ |\ a\in A\right\} $}
בתור \L{$U\times V$} אם \L{$A$} כוללת יותר מאיבר אחד(, אבל אנחנו
כן אומרים שכל קבוצה פתוחה יכולה להתקבל על ידי איחוד אינסופי של קבוצות
בסיס. למשל, ב-\L{$\mathbb{R}^{2}$} כל איבר מהצורה \L{$U\times V$}
הוא \textbf{מלבן} )מלבן פתוח; כלומר, בלי השפה שלו( אבל אם נאחד אינסוף
מלבנים שייבחרו בצורה זהירה נוכל לקבל עיגול פתוח.

כזכור, אי אפשר סתם לקרוא בשם \textquotedblleft בסיס\textquotedblright{}
לכל קבוצה, יש שתי דרישות שצריכות להתקיים. אחת היא שכל איבר במרחב יהיה
שייך לאחת מהקבוצות וזה כמובן עובד כאן כי \L{$X\times Y$} עצמה היא
איבר בסיס. הדרישה השניה טיפה יותר טריקית: אם \L{$B_{1},B_{2}$} הם
אברי בסיס ו-\L{$x\in B_{1}\cap B_{2}$} אז צריך להראות שקיים איבר
בסיס \L{$B$} כך ש-\L{$x\in B\subseteq B_{1}\cap B_{2}$}. במקרה שלנו
זה עובד כי \L{$B_{1}\cap B_{2}$} הוא די פשוט עבור קבוצות שהן מכפלה
- זה שוב פעם תורת קבוצות בסיסית בפעולה:

\L{$\left(U_{1}\times V_{1}\right)\cap\left(U_{2}\times V_{2}\right)=\left(U_{1}\cap U_{2}\right)\times\left(V_{1}\cap V_{2}\right)$}

לכן תמיד אפשר לבחור \L{$B=B_{1}\cap B_{2}$} כי גם הוא איבר בסיס,
וזה מסיים את הסיפור הזה.

מה קורה אם לוקחים את \L{$\mathbb{R}$} עם הטופולוגיה הרגילה עליו )טופולוגיה
שאנחנו כבר יודעים לקבל בעזרת מטריקות ובעזרת טופולוגיית הסדר( ומסתכלים
על הטופולוגיה של המכפלה \L{$\mathbb{R}^{n}$} שמתקבלת ממה שתיארנו
כאן? באופן לא מפתיע, מקבלים את הטופולוגיה ה\textquotedblright רגילה\textquotedblright{}
על \L{$\mathbb{R}^{n}$}, זו שמתקבלת מהמטריקה האוקלידית. אז אין לנו
הזיות כמו מה שקרה קודם עם טופולוגיית תת-המרחב שבו הבניה החדשה שלנו
התנגשה חזיתית עם בניות קודמות. בינתיים.

יש עוד דרך מועילה מאוד במתמטיקה באופן כללי להסתכל על מכפלות - להסתכל
על \textbf{הטלות}. אם יש לנו מרחב מכפלה \L{$X\times Y$} אז אנחנו
מגדירים הטלות ברכיב הראשון והשני על ידי
\begin{itemize}
\item \L{$\pi_{1}:X\times Y\to X$} כך ש-\L{$\pi_{1}\left(x,y\right)=x$}
\item \L{$\pi_{2}:X\times Y\to Y$} כך ש-\L{$\pi_{2}\left(x,y\right)=y$}
\end{itemize}
את זה מן הסתם קל להכליל גם למכפלה של מספר סופי כלשהו של מוכפלים:
\begin{itemize}
\item \L{$\pi_{i}:X_{1}\times X_{2}\times\ldots\times X_{n}\to X_{i}$}
כך ש-\L{$\pi_{i}\left(x_{1},\ldots,x_{n}\right)=x_{i}$}
\end{itemize}
עכשיו, נניח שעדיין אין לנו טופולוגיה על \L{$X\times Y$}, אבל יש לנו
יעד קונקרטי - אנחנו רוצים שההטלות מהמכפלה הזו יהיו \textbf{פונקציות
רציפות}. כלומר, אני רוצה ש-\L{$\pi^{-1}_{1}\left(U\right)$} תהיה
קבוצה פתוחה ב-\L{$X\times Y$} לכל \L{$U\subseteq X$} פתוחה, ובדומה
גם עבור \L{$\pi_{2}$}. עכשיו, מה זו \L{$\pi^{-1}_{1}\left(U\right)$}?
קל לבדוק ש-\L{$\pi^{-1}_{1}\left(U\right)=U\times Y$}, אז אני רוצה
שכל הקבוצות מהצורה הזו יהיו פתוחות. בדומה אני רוצה שכל הקבוצות מהצורה
\L{$X\times V$}, כאשר \L{$V$} פתוחה ב-\L{$Y$}, יהיו פתוחות. הגדרה
כזו, של \textquotedblleft הנה אוסף קבוצות שאני רוצה שיהיו פתוחות,
תנו לי את הטופולוגיה המינימלית שמקיימת את זה\textquotedblright{} היא
פשוט הגדרה של טופולוגיה בעזרת \textbf{תת-בסיס}. הדרישה היחידה מתת-בסיס
הוא שכל איבר במרחב יהיה שייך אל איבר כלשהו ממנו מה שמן הסתם מתקיים
עבור \L{$\pi^{-1}_{1}\left(X\right)=X\times Y$}. עכשיו, מכיוון שכל
האיברים \L{$U\times Y$} ו-\L{$X\times V$} כבר היו שייכים לבסיס שתיארתי
קודם, והחיתוך של שני איברים כאלו הוא משהו מהצורה \L{$U\times V$},
אני מקבל מההגדרה הזו \textbf{בדיוק} את אותה טופולוגיה כמו קודם. פשוט
מאוד לבינתיים!

אוקיי, אז בואו נסבך את זה עם האופן שבו מתמטיקאים אוהבים לסבך דברים:
בואו נכניס לתמונה את \textbf{אינסוף}.

\section{טופולוגיית המכפלה - החלק האינסופי}

האופן שבו הגדרתי מכפלה קרטזית, עם זוגות סדורים ואז \textquotedblleft בואו
נעשה את זה שוב\textquotedblright{} לא מביא אותנו יותר מדי רחוק, פורמלית.
אפילו הקטע של ה\textquotedblright בואו נעשה את זה שוב\textquotedblright .
הוא בעייתי. \href{https://gadial.net/2020/06/20/general_cartesian_products/}{יש לי פוסט על זה},
אז כאן אני אסתפק בשורה התחתונה: אם רוצים להגדיר מכפלה קרטזית של יותר
משתי קבוצות, הדרך הנכונה לעשות את זה היא באמצעות \textbf{פונקציות}
)ופונקציות, בתורן, הן מושג שמסתמך על מכפלה קרטזית של שתי קבוצות אז
גם למושג הזה אנחנו זקוקים(. פורמלית, נניח שיש לנו קבוצת \textbf{אינדקסים}
\L{$\Lambda$} כך שלכל \L{$\alpha\in\Lambda$} מתאימה קבוצה \L{$X_{\alpha}$}.
אז \textbf{המכפלה הקרטזית} של הקבוצות הללו, שמסומנת \L{$\prod_{\alpha\in\Lambda}X_{\alpha}$}
היא אוסף כל הפונקציות \L{$a:\Lambda\to\bigcup_{\alpha\in\Lambda}X_{\alpha}$}
כך ש-\L{$a\left(\alpha\right)\in X_{\alpha}$} לכל \L{$\alpha\in\Lambda$}.
הרעיון הוא ש-\L{$\Lambda$} מייצגת את ה\textbf{קואורדינטות} של המכפלה
ו-\L{$a$} היא מה שאומר לנו, לכל קואורדינטה, איזה ערך ספציפית נמצא
בה - תחת הדרישה שהערך בקואורדינטה שמאונדקסת על ידי \L{$\alpha$} שייך
לקבוצה \L{$X_{\alpha}$}. זה מכליל את מה שראינו קודם: המקרה של \L{$X_{1}\times\ldots\times X_{n}$}
הוא פשוט מה שמקבלים עבור \L{$\Lambda=\left\{ 1,2,\ldots,n\right\} $},
אבל אפשר לקחת את זה הרבה יותר רחוק - לקבוצות אינדקסים אינסופיות, ואפילו
לא בנות מניה.

הדוגמא האינסופית הפשוטה ביותר היא \L{$X_{1}\times X_{2}\times X_{3}\times\ldots$}
שמתאימה למקרה של \L{$\Lambda=\mathbb{N}$} )בגרסה של \L{$\mathbb{N}$}
שלא כוללת את {\beginL 0\endL}(. למשל, אם ניקח \L{$X_{1}=X_{2}=X_{3}=\ldots=\mathbb{R}$}
אז למכפלה שמקבלים קוראים \L{$\mathbb{R}^{\omega}$} )פורמלית, \L{$\omega=\left\{ 0,1,2,\ldots\right\} $}
אבל נעזוב את זה(. עוד נשתמש במרחב הזה אבל אותי יותר מעניין להראות
דוגמא פחות מובנת מאליה במבט ראשון: בואו ניקח את \L{$\mathbb{R}$}
ונכפול אותו בעצמו מספר לא בן מניה של פעמים, עם קבוצת האינדקסים \L{$\Lambda=\mathbb{R}$}.
מה מקבלים? ובכן, מקבלים את כל הפונקציות מ-\L{$\mathbb{R}$} )זו קבוצת
האינדקסים( אל \L{$\mathbb{R}$} )זה האיחוד של כל העותקים של \L{$\mathbb{R}$}
שאנחנו מכפילים( כך שלכל \L{$r\in\mathbb{R}$} מתקיים \L{$a\left(r\right)\in\mathbb{R}$}
)כלומר, במקרה הזה הדרישה על שייכות ל-\L{$X_{\alpha}$} היא טריוויאלית
כי כל המוכפלים זהים(. כלומר, קיבלנו... פונקציות רגילות על \L{$\mathbb{R}$}?
את הקבוצה \L{$\mathbb{R}^{\mathbb{R}}$}? אלו הפונקציות שמככבות בחשבון
אינפיניטסימלי {\beginL 1\endL}, נקודת המוצא שלנו לכל המהומה הנוכחית,
ואפשר לחשוב עליהן גם בתור איברים של מכפלה קרטזית של \L{$\mathbb{R}$}
בעצמו מספיק פעמים.

יפה, אז אנחנו מבינים איך להגדיר מכפלה קרטזית כללית. עכשיו, בהינתן
מכפלה כזו \L{$\prod_{\alpha\in\Lambda}X_{\alpha}$} , איך נגדיר עליה
טופולוגיה? בחלק הקודם הצגתי שתי גישות שנתנו את אותה התוצאה:
\begin{enumerate}
\item לקחת בתור בסיס לטופולוגיה את כל הקבוצות שהן מכפלות של קבוצות פתוחות.
כלומר, כל הקבוצות מהצורה \L{$\prod_{\alpha\in\Lambda}U_{\alpha}$}
כך ש-\L{$U_{\alpha}\subseteq X_{\alpha}$} פתוחה.
\item לקחת בתור תת-בסיס לטופולוגיה את כל הקבוצות שהן מקורות של פונקציות
ההטלה. כלומר, כל הקבוצות מהצורה \L{$\pi^{-1}_{\alpha}\left(U_{\alpha}\right)$}
כך ש-\L{$U_{\alpha}\subseteq X_{\alpha}$} פתוחה.
\end{enumerate}
רק כדי ליישר קו, הטלה מוגדרת כמו קודם: אם \L{$a\in\prod_{\alpha\in\Lambda}X_{\alpha}$}
הוא איבר במכפלה הקרטזית, אז \L{$\pi_{\alpha}\left(f\right)=a\left(\alpha\right)$}.
אז אני מקווה ששתי ההגדרות ברורות - והנה העניין, \textbf{הן לא נותנות
את אותה טופולוגיה}.
\begin{enumerate}
\item ההגדרה הראשונה נותנת את מה שנקרא \textbf{טופולוגיית התיבה} )\L{Box
Topology}(
\item ההגדרה השניה נותנת את מה שנקרא \textbf{טופולוגיית המכפלה} )\L{Product
Topology}(
\end{enumerate}
מה ההבדל? ובכן, בואו נראה מה קורה בטופולוגיית המכפלה. כל איבר של תת-הבסיס
שלנו הוא מהצורה \L{$\pi^{-1}_{\alpha}\left(U_{\alpha}\right)$}. אינטואיטיבית
מה שזה אומר הוא שהמקור הזה הוא מכפלה של \textbf{כל} המרחב עבור \textbf{כל}
הקואורדינטות, חוץ מאשר עבור הקואורדינטה ה-\L{$\alpha$} ושם מי שמוכפל
הוא \L{$U_{\alpha}$}. למשל, עבור המכפלה \L{$X_{1}\times X_{2}\times X_{3}\times\ldots$}
אנחנו נקבל ש-

\L{$\pi^{-1}_{2}\left(U\right)=X_{1}\times U\times X_{3}\times\ldots$}

כלומר, אנחנו לוקחים את כל המרחב ו\textbf{מגבילים} רק קואורדינטה אחת.
לעומת זאת, בטופולוגיית התיבה אנחנו יכולים להגביל את כל אינסוף הקואורדינטות
\textbf{בו זמנית}.

עכשיו, בהגדרה השניה הגדרנו \textbf{תת-בסיס}. זה אומר שעדיין אפשר לקחת
חיתוכים של האיברים הללו. כשאנחנו חותכים שני איברים כאלו אנחנו נקבל
הגבלה על \textbf{שתי} קואורדינטות, ובאופן כללי חיתוך סופי של איברים
כאלו יתן לנו הגבלה על \textbf{מספר סופי} של קואורדינטות; אבל בטופולוגיית
התיבה אפשר להגביל גם \textbf{מספר אינסופי} של קואורדינטות. זה ההבדל
בין שתי ההגדרות, ועכשיו גם ברור למה במקרה הסופי הוא לא בא לידי ביטוי
בכלל.

לאלו מכן שמכירות אלגברה זה אולי מזכיר את ההבדל בין\textbf{ מכפלה ישרה}
ובין \textbf{סכום ישר} של מבנים אלגבריים; מכפלה ישרה זו פשוט מכפלה
קרטזית רגילה שיורשת את המבנה האלגברי שלה מהאיברים שמשתתפים במכפלה,
אבל \textbf{סכום ישר} כולל רק את האיברים של המכפלה הישרה שהם {\beginL 0\endL}
בכל כניסה חוץ ממספר סופי של כניסות. עוד מקרה שבו ההבדל הזה נפוץ ומוכר
הוא ההבדל שבין \textbf{פולינומים} ובין \textbf{טורי חזקות}; אפשר לחשוב
על פולינום בתור טור חזקות שבו כל המקדמים הם {\beginL 0\endL} החל ממקום
מסוים. בקיצור, הרעיון הזה הוא לא כזה חדש ומפתיע - אבל במבט ראשון זה
כן מוזר שיש פה הבדל, ועוד יותר לא ברור למה להעדיף דווקא את ההגדרה
השניה, הרי הראשונה מרגישה טבעית יותר. הנימוק הוא כמובן \textquotedblleft חכו
ותראו, משפט טיכונוף לא עובד בטופולוגיית התיבה!!!!!\textquotedblright{}
)ההגדרה של טופולוגיית המכפלה ליטרלי הגיעה מהמאמר של טיכונוף שבו הוא
הוכיח את משפט טיכונוף( אבל מכיוון שאין לנו מושג כרגע מה זה משפט טיכונוף
אולי כדאי לנסות לתת הסבר קונקרטי יותר כבר עכשיו.

ראשית, מה ההבדל בין הטופולוגיות? אפשר להציג את שניהן באמצעות בסיסים.
בטופולוגיית התיבה, כל איבר בסיס הוא מהצורה \L{$\prod_{\alpha\in\Lambda}U_{\alpha}$}
כך ש-\L{$U_{\alpha}\subseteq X_{\alpha}$} פתוחה. בטופולוגיית המכפלה,
כל איבר בסיס הוא מהצורה \L{$\prod_{\alpha\in\Lambda}U_{\alpha}$}
כך ש-\L{$U_{\alpha}\subseteq X_{\alpha}$} פתוחה ובנוסף לכך מתקיים
\L{$U_{\alpha}=X_{\alpha}$} לכל \L{$\alpha\in\Lambda$} פרט אולי
למספר סופי. מה שברור מייד מלהציג את זה ככה הוא שכל קבוצה שפתוח בטופולוגיית
המכפלה פתוחה גם בטופולוגיית התיבה. אינטואיטיבית אולי היינו משתמשים
במילים יפות כמו \textquotedblleft טופולוגיית התיבה \textbf{עשירה יותר}
מטופולוגיית המכפלה\textquotedblright{} או \textquotedblleft טופולוגיית
התיבה \textbf{עדינה יותר }מטופולוגיית המכפלה\textquotedblright{} אבל
אני גם לא סובל את השמות הללו וגם לפעמים משתמשים בהם במשמעות ההפוכה.
עדיין, לא ברור למה טופולוגיית המכפלה עדיפה - האם לא עדיפה טופולוגיית
התיבה, שבה יש יותר קבוצות פתוחות? יותר קבוצות פתוחות, חיים יותר מעניינים,
לא?

אז ראשית, זה לא שיש טופולוגיה טובה וטופולוגיה רעה. לכל אחת מהן יש
שימושים. פשוט אחת מהן הרבה יותר שימושית מהאחרת. אז בואו נדגים משפט
שנראה \textquotedblleft טבעי\textquotedblright{} ואפשר להוכיח אותו
בטופולוגיית המכפלה אבל אי אפשר להוכיח בטופולוגיית התיבה, ודווקא מקבלים
אינטואיציה טובה למה \textquotedblleft עוד קבוצות פתוחות\textquotedblright{}
זה לאו דווקא דבר טוב.

העניין הוא כזה: איך אפשר לבנות פונקציה רציפה \textbf{שהטווח} שלה הוא
מכפלה קרטזית? למשל, \L{$f:\mathbb{R}\to\mathbb{R}^{2}$}? תוצאה בסיסית
בחשבון אינפיניטסימלי אומרת שכדאי להסתכל על \textbf{הרכיבים} של הפונקציה
- להציג אותה בתור \L{$f\left(x\right)=\left(f_{1}\left(x\right),f_{2}\left(x\right)\right)$}
כך ש-\L{$f_{1},f_{2}:\mathbb{R}\to\mathbb{R}$}, ואם כל אחת מהפונקציות
הללו היא רציפה בעצמה, אז \L{$f$} המקורית רציפה. בואו נוכיח את זה
גם למקרה שבו יש לנו מכפלה קרטזית אינסופית. כלומר נתונן לנו מרחב טופולוגי
\L{$X$} כלשהו ומשפחה \L{$Y_{\alpha}$} של מרחבים ואנחנו מסתכלים על
פונקציה \L{$f:X\to\prod_{\alpha\in\Lambda}Y_{\alpha}$}. את הפונקציה
הזו נתאר על ידי רכיבים, כך ש-\L{$f_{\beta}:X\to Y_{\beta}$} הוא הרכיב
שמתאים ל-\L{$\beta$}. פורמלית, \L{$f_{\beta}=\pi_{\beta}f$} - ההרכבה
של ההטלה על הרכיב \L{$\beta$} שמורכבת על הפונקציה \L{$f$}. שימו
לב שאם \L{$f$} רציפה ואנחנו מעל טופולוגיית המכפלה ולכן גם \L{$\pi_{\beta}$},
אז נקבל מייד ש-\L{$f_{\beta}$} רציפה כי היא הרכבה של שתי פונקציות
רציפות והראיתי בפוסט הקודם שהרכבה כזו היא רציפה. החלק המעניין הוא
הכיוון השני - נניח שכל הרכיבים הללו רציפים ונוכיח ש-\L{$f$} רציפה.

בשביל להראות ש-\L{$f$} היא רציפה, ניקח קבוצה פתוחה ב-\L{$\prod_{\alpha\in\Lambda}Y_{\alpha}$}
ונראה שהמקור שלה גם פתוח. הנה מייד ברור מה היתרון בטופולוגיית המכפלה
- אנחנו צריכים לטפל רק בקבוצות \textbf{פשוטות יחסית}. בטופולוגיית
התיבה האתגר שלנו יהיה גדול בהרבה דווקא בגלל שיש יותר קבוצות פתוחות.
עכשיו, בפוסט הקודם הראיתי שלא צריך לטפל \textbf{בכל} קבוצה פתוחה;
אם יש למרחב שלנו תת-בסיס \L{$\mathcal{B}$}, מספיק להראות שלכל \L{$B\in\mathcal{B}$},
\L{$f^{-1}\left(B\right)$} היא קבוצה פתוחה. עכשיו, בטופולוגיית המכפלה
תת-הבסיס שלנו מורכב מקבוצות מהצורה \L{$\pi^{-1}_{\beta}\left(U\right)$}
כש-\L{$U\subseteq Y_{\beta}$} היא קבוצה פתוחה. עכשיו, עוד דבר שראינו
בפוסט הקודם הוא שמתקיים באופן כללי, עבור פונקציות \L{$f:A\to B$}
ו-\L{$g:B\to C$} ו-\L{$D\subseteq C$}:

\L{$\left(gf\right)^{-1}\left(D\right)=f^{-1}\left(g^{-1}\left(D\right)\right)$}

אצלנו \L{$g$} היא הפונקציה \L{$\pi_{\beta}:\prod_{\alpha\in\Lambda}Y_{\alpha}\to Y_{\beta}$}.
מה ההרכבה \L{$\pi_{\beta}f$}? כאמור, זו בדיוק פונקציית הרכיב \L{$f_{\beta}$}
שהנחתי שהיא רציפה, ומכאן ש-\L{$f^{-1}_{\beta}\left(U\right)$} קבוצה
פתוחה, ולכן \L{$f^{-1}\left(\pi^{-1}_{\beta}\left(U\right)\right)=f^{-1}_{\beta}\left(U\right)$}
היא קבוצה פתוחה וזה בדיוק מה שרצינו להראות - שהמקור על פי \L{$f$}
של איבר תת-הבסיס \L{$\pi^{-1}_{\beta}\left(U\right)$} הוא קבוצה פתוחה.
זה מסיים את ההוכחה הזו.

ההוכחה הזו אולי נתנה לנו אינטואיציה לקושי הנוסף שטופולוגיית התיבה
מציבה בפנינו, אבל זה לכשעצמו לא \textbf{מוכיח} שהטענה פשוט לא עובדת
עבור טופולוגיית התיבה. צריך לתת דוגמא מפורשת, והדוגמא הזו תצטרך מן
הסתם להיות מכפלה אינסופית )אחרת שתי הטופולוגיות שקולות(. אז בואו ניקח
את המקרה הפשוט של \L{$\mathbb{R}^{\omega}$} שבו \L{$\mathbb{R}$}
מוכפל בעצמו מספר בן מניה של פעמים. עכשיו, את מי ניקח בתור דוגמא לפונקציה
שהרכיבים שלה רציפים אבל היא עצמה לא רציפה? ובכן, זוכרים שקודם השתמשתי
ב-\L{$\left\{ \left(a,a\right)\ |\ a\in A\right\} $} בתור דוגמא לתת-קבוצה
של מכפלה שאי אפשר \textquotedblleft לפרק למכפלה\textquotedblright ?
נעשה משהו דומה אבל עם פונקציה - נגדיר

\L{$f:\mathbb{R}\to\mathbb{R}^{\omega}$}

על ידי

\L{$f\left(x\right)=\left(x,x,x,\ldots\right)$}

כל רכיב כאן הוא פונקציית הזהות \L{$f_{i}\left(x\right)=x$} שהיא בוודאי
רציפה על \L{$\mathbb{R}$}. אז למה \L{$f$} לא רציפה בטופולוגיית התיבה?
אני חושב שקל להבין את זה אם מסתכלים רגע באופן נקדותי על הפונקציה בנקודה
כלשהית וספציפית נוח להסתכל ב-{\beginL 0\endL} בגלל הסימטריה. אם \L{$f$}
רציפה ב-{\beginL 0\endL} זה אומר שלכל קבוצה פתוחה \L{$V\subseteq\mathbb{R}^{\omega}$}
שמכילה את \L{$\left(0,0,0,\ldots\right)$}, קיימת קבוצה פתוחה \L{$U\subseteq\mathbb{R}$}
שמכילה את {\beginL 0\endL} כך ש-\L{$f\left(x\right)\in V$} לכל \L{$x\in U$},
כלומר \L{$\left(x,x,x,\ldots\right)\in V$} לכל \L{$x\in U$}. בזכות
העובדה ש-\L{$U$} פתוחה ב-\L{$\mathbb{R}$}, יש בה \L{$x\ne0$} כלשהו,
ולכן אנחנו מצטמצמים לשאלה הבאה: האם נוכל להטריל את נסיון ההוכחה הזה
על ידי כך שנמצא קבוצה \L{$V$} שמכילה את \L{$\left(0,0,0,\ldots\right)$}
אבל לא מכילה \textbf{שום איבר אחר} מהצורה \L{$\left(x,x,x,\ldots\right)$}?
והתשובה היא \textquotedblleft כן, בקלות\textquotedblright{} כי פשוט
צריך לבנות את \L{$V$} בתור אוסף של קטעים פתוחים סביב {\beginL 0\endL}
שהולכים ומצטמצמים. פורמלית:

\L{$V=\left(-1,1\right)\times\left(-\frac{1}{2},\frac{1}{2}\right)\times\left(-\frac{1}{3},\frac{1}{3}\right)\times\dots\times\left(-\frac{1}{n},\frac{1}{n}\right)\times\dots$}

הסיבה שאנחנו מסוגלים לבנות דבר כזה היא שטופולוגיית התיבה נותנת לנו
לשלוט על כל אינסוף הרכיבים במכפלה הזו. בטופולוגיית התיבה זה בלתי אפשרי
- אחרי מספר סופי של רכיבים נצטרך להתייאש וכל יתר הרכיבים יהיו פשוט
\L{$\mathbb{R}$} ונסיון ההטרלה ייכשל.

לסיום, בואו נראה עוד דרך להסתכל על זה, הפעם עבור המרחב \L{$\mathbb{R}^{\mathbb{R}}$}
שהוא כזכור מצד אחד פשוט אוסף הפונקציות הממשיות ומצד שני אנחנו רואים
אותו בתור מרחב מכפלה עם טופולוגיה עליו. מה זו קבוצה פתוחה במרחב הזה?
בטופולוגיית המכפלה, קבוצת תת-בסיס היא \L{$\mathbb{R}$} בכל הקואורדינטות
חוץ מאחת, שבה היא \L{$U$} עבור קבוצה פתוחה \L{$U\subseteq\mathbb{R}$}.
נאמר שזה בקואורדינטה \L{$a$}. אז \L{$f\in\mathbb{R}^{\mathbb{R}}$}
שייכת לאיבר תת-הבסיס הזה אם היא מקיימת \L{$f\left(a\right)\in U$}.
כלומר, השייכות לקבוצה הפתוחה הזו מתורגמת לאילוץ \textbf{נקודתי}: הערך
של \L{$f$} בנקודה ספציפית לא יכול להיקבע בחופשיות אלא הוא שייך ל-\L{$U$}.
אם נחשוב על \L{$U$} בתור אוסף של קטעים פתוחים, אז נוח לי לדמיין את
הקבוצה הפתוחה בתור מעין \textquotedblleft מחסום\textquotedblright{}
אנכי שהוא הישר \L{$x=a$}, ובתוך המחסום הזה יש \textquotedblleft שערים\textquotedblright{}
שדרכם הפונקציה יכולה לעבור ומתאימים לקטעים הפתוחים ב-\L{$U$}. 

כשלוקחים אברי תת-בסיס כזה וחותכים אותם, זה כאילו מוספים עוד מספר סופי
של מחסומים עם שערים שכאלו. ובטופולוגיית התיבה יכולים להיות מחסומים
\textbf{בכל מקום}, מה שמאפשר לנו להגביל מאוד את הפונקציות שיכולות
להיכנס לקבוצות הפתוחות שאנחנו בונים.

מעניין לראות איך זה בא לידי ביטוי בהקשר של \textbf{התכנסות}. המרחב
שלנו הוא מרחב של פונקציות ממשיות, אז כשאני מדבר על התכנסות אני מדבר
על לקחת סדרה של פונקציות ולבדוק לאיזו \textbf{פונקציה} היא מתכנסת.
זה משהו \href{https://gadial.net/2024/02/24/sequences_and_series_of_functions/}{שכבר דיברתי עליו בבלוג},
ואמרתי אז שיש שני סוגי התכנסות שמעניינים אותנו בהקשר הזה; שתי דרכים
לומר שהסדרה \L{$f_{n}$} של פונקציות ממשיות מתכנסת אל הפונקציה \L{$f$}:
\begin{itemize}
\item \L{$f_{n}$} מתכנסת \textbf{נקודתית} אל \L{$f$} אם לכל \L{$x\in\mathbb{R}$},
לכל \L{$\varepsilon>0$} קיים \L{$N$} טבעי כך שלכל \L{$n>N$} מתקיים
\L{$\left|f_{n}\left(x\right)-f\left(x\right)\right|<\varepsilon$}.
\item \L{$f_{n}$} מתכנסת \textbf{במידה שווה} אל \L{$f$} אם לכל \L{$\varepsilon>0$}
קיים \L{$N$} טבעי כך שלכל \L{$n>N$} מתקיים \L{$\left|f_{n}\left(x\right)-f\left(x\right)\right|<\varepsilon$}
לכל \L{$x\in\mathbb{R}$}.
\end{itemize}
המשמעות של התכנסות בטופולוגיית המכפלה היא התכנסות \textbf{נקודתית}.
למה? כי אם \L{$f_{n}\to f$} בטופולוגיית המכפלה, אז לכל סביבה של \L{$f$}
קיים \L{$N$} כך שלכל \L{$n>N$} גם \L{$f_{n}$} שייך לאותה סביבה.
עכשיו, נניח שאני רוצה להראות ש-\L{$f_{n}$} מתכנסת נקודתית, אז ניקח
נקודה \L{$x\in\mathbb{R}$} ו-\L{$\varepsilon>0$} ועכשיו ניקח את
הסביבה שהיא איבר תת-בסיס שהוא \L{$\mathbb{R}$} בכל קואורדינטה למעט
זו של \L{$x$} שבה יש את \L{$\left(f\left(x\right)-\varepsilon,f\left(x\right)+\varepsilon\right)$}.
אז קיים \L{$N$} כך שלכל \L{$n>N$} מתקים \L{$f_{n}\left(x\right)\in\left(f\left(x\right)-\varepsilon,f\left(x\right)+\varepsilon\right)$}
שזו דרך אחרת לומר \L{$\left|f_{n}\left(x\right)-f\left(x\right)\right|<\varepsilon$}.

בכיוון השני, אם \L{$f_{n}\to f$} נקודתית זה מספיק כדי להבטיח התכנסות
בטופולוגיית המכפלה. אם לוקחים סביבה של \L{$f$}, אז יש איבר בסיס שמכיל
את \L{$f$} ומוכל בסביבה הזו. איבר הבסיס הזה יכול לכלול הגבלות לא
רק על קואורדינטה אחת אלא על מספר \textbf{סופי} של קואורדינטות, אבל
יש טריק סטנדרטי של אינפי להתמודד עם זה. לכל קואורדינטה יש את ה-\L{$N_{x}$}
משלה, על פי ההתכנסות הנקודתית; פשוט נבחר את \L{$N=\max_{x}\left\{ N_{x}\right\} $}
להיות המקסימום על כל ה-\L{$N_{x}$}-ים הללו. המקסימום קיים כי יש רק
מספר \textbf{סופי} של \L{$N_{x}$}-ים שכאלו.

כל זה כמובן לא יכול לעבוד בטופולוגיית התיבה. זה השלב שבו היה מאוד
נחמד אם הייתי אומר שמה שטופולוגיית התיבה נותנת לנו הוא התכנסות במידה
שווה - אבל לא, היא לא נותנת אפילו את זה. גם התכנסות במידה שווה היא
חלשה מדי מכדי להבטיח התכנסות בטופולוגיית התיבה )יש טופולוגיה אחרת,
\L{uniform topology}, שבאה לתת התכנסות במידה שווה, אבל לא נדבר עליה
כאן(.

למה זה לא עובד אפילו עם התכנסות במידה שווה? ובכן, אותו טריק של

\L{$V=\left(-1,1\right)\times\left(-\frac{1}{2},\frac{1}{2}\right)\times\left(-\frac{1}{3},\frac{1}{3}\right)\times\dots\times\left(-\frac{1}{n},\frac{1}{n}\right)\times\dots$} 

שראינו קודם יעבוד גם כאן. פשוט במקום עכשיו שזו תהיה \textbf{כל} \L{$V$},
זו תהיה \L{$V$} רק בכניסות שהם מספרים טבעיים. כלומר, אנחנו נדרוש
ש-\L{$f_{n}$}, בקואורדינטה \L{$k$} עבור \L{$k$} טבעי, תהיה קרובה
אל \L{$f$} עד כדי \L{$\frac{1}{k}$} )\L{$\left|f_{n}\left(k\right)-f\left(k\right)\right|<\frac{1}{k}$}(.
התכנסות במידה שווה לא יכולה לתת את זה - היא צריכה \L{$\varepsilon$}
קונקרטי שיעבוד בו זמנית לכל הקואורדינטות, אבל כאן צריך שככל שהאינדקס
של הקואורדינטה גדל, כך ה-\L{$\varepsilon$} הרלוונטי אליה יהיה יותר
ויותר קטן. במילים אחרות, צריך שהוא יהיה \L{$\varepsilon\left(x\right)$}
, כזה שתלוי ב-\L{$x$} ויכול לקטון כרצוננו. אפשר לנסח התכנסות כזו,
אבל אני חושב שכבר ברור איך טופולוגיית התיבה מסבכת לנו את החיים ולמה
בדרך כלל נעדיף להשתמש בטופולוגיית המכפלה.
\end{document}
